\begin{table}[t]
\centering
\caption{The dollar-boundary effect under five ways of handling the roundness of the price a step lands on.}
\label{tab:specs}
\small
\begin{tabular}{p{6.0cm}lrc}
\toprule
Specification & Roundness assumption & Estimate & 95\% CI \\
\midrule
Matched design (E0) & none & -6.14* & [-8.36, -3.86] \\
Step 10 minus the mean placebo step & $.00$ weakly most attractive & -2.50* & [-3.91, -1.28] \\
Ordinal roundness score & linear in $\rho$ & -3.25* & [-5.64, -0.99] \\
Measured corpus frequency of the destination & linear in log frequency & -3.52* & [-5.44, -1.73] \\
Step 10 minus the largest placebo step & none (not a bound) & -3.48* & [-5.36, -2.32] \\
\bottomrule
\end{tabular}
\vspace{2pt}
\parbox{\linewidth}{\footnotesize Row 1 uses no roundness model: it compares two step-10 jumps that both end on $.00$ and differ only in whether the leftmost digit changes, estimated on E0's own data. Rows 2 to 5 use the core sweep. Row 2 is the uncorrected contrast. If a whole-dollar destination is at least as attractive as every within-dollar destination, its expectation lies at or above the true effect, so it can only understate a negative effect. Row 3 regresses the nine placebo jumps on the ordinal roundness score and extrapolates to a whole dollar; the fit rests on eight placebos with $\rho = 1$ and one with $\rho = 2$. Row 4 replaces the score with the measured corpus frequency of each destination, giving nine distinct values. Row 5 subtracts the largest of the nine placebo jumps; because the largest of nine noisy jumps tends to exceed its expectation, this contrast tends to overstate the effect and is shown for completeness only. Rows 2 to 5 are bootstrapped on the same 400 block resamples.}
\end{table}
